\documentclass[11pt,a4paper]{article}

\usepackage{amsmath,amssymb,amsthm}
\usepackage[margin=2.5cm]{geometry}
\usepackage{enumitem}
\usepackage{booktabs}

\newcommand{\dd}{\mathrm{d}}
\newtheorem{proposition}{Proposition}

\title{Right-Truncation Biases the Excess Variance Downwards\\[4pt]
\large Why dropping long spells manufactures positive duration dependence}
\author{}
\date{}

\begin{document}
\maketitle

\section{Statement of the problem}

Our decomposition requires \emph{two completed} unstable spells per
individual, observed inside a finite window of length~$T$.  Any individual
whose spell runs past the end of the window contributes no completed spell
and drops out of the estimation sample.  The observed duration distribution
is therefore not the population distribution but its
\textbf{right-truncation} at~$T$.

This note shows that such truncation biases the excess-variance component
$\sigma^2_e$ \emph{downwards, always}, and therefore manufactures apparent
\emph{positive} duration dependence (a rising hazard) even when the true
process has none.

\paragraph{Truncation, not censoring.}
The distinction matters.  Under \emph{censoring} the observation is retained
and flagged as incomplete, and a correctly specified likelihood recovers the
truth.  Under \emph{truncation} the observation is discarded outright, and
no amount of care with the remaining data undoes the selection.  The
two-completed-spell requirement is a truncation mechanism.

\paragraph{Sign convention.}
Throughout, $\sigma^2_e > 0$ denotes \emph{negative} duration dependence
(hazard falling with duration, over-dispersion relative to the exponential),
and $\sigma^2_e < 0$ denotes \emph{positive} duration dependence (hazard
rising, under-dispersion).  The bias established below is towards
$\sigma^2_e < 0$.

\section{What $\sigma^2_e$ measures}

Let individual~$i$ contribute two spells $y_{i1}, y_{i2}$ drawn i.i.d.\ from
$F_i$ with mean $\mu_i$ and variance $\sigma^2_i$.  With two observations the
natural estimators are
\[
  \hat\mu_i = \frac{y_{i1}+y_{i2}}{2},
  \qquad
  \hat\sigma^2_i = \frac{(y_{i1}-y_{i2})^2}{2},
\]
whose expectations are
\begin{equation}
  E[\hat\sigma^2_i] = \sigma^2_i,
  \qquad
  E[\hat\mu_i^{\,2}] = \mu_i^2 + \tfrac{1}{2}\sigma^2_i .
  \label{eq:moments}
\end{equation}
The excess-variance estimator is
$\hat\sigma^2_e = \frac{1}{n}\sum_i \bigl(\tfrac{3}{2}\hat\sigma^2_i - \hat\mu_i^{\,2}\bigr)$,
so by \eqref{eq:moments}
\begin{equation}
  \boxed{\;
  \sigma^2_e \;=\; E_i\!\left[\sigma^2_i - \mu_i^2\right]
  \;=\; E_i\!\left[\mu_i^2\left(CV_i^2 - 1\right)\right].
  \;}
  \label{eq:sigmae}
\end{equation}
Likewise
$\hat\sigma^2_c = \frac{1}{n}\sum_i\bigl(\hat\mu_i^{\,2}-\tfrac12\hat\sigma^2_i\bigr)$
has expectation $E_i[\mu_i^2]$, the within-variance that would obtain if
every individual were exponential.

Equation~\eqref{eq:sigmae} is the whole content of the diagnostic: in levels,
$\sigma^2_e$ is a test of $CV_i = 1$ against the exponential benchmark,
averaged over individuals.  Anything that compresses individual dispersion
relative to the individual mean drives it negative.

\section{The mechanism in one line}

For any distribution with density $f$ and survival $S$, truncated at $T$, the
hazard of the \emph{observed} variable is
\begin{equation}
  h_T(y) \;=\; \frac{f(y)}{S(y) - S(T)},
  \qquad y < T,
  \qquad\text{so}\qquad
  h_T(y) \;\xrightarrow[y \to T]{}\; \infty .
  \label{eq:hazard}
\end{equation}
Truncation forces every surviving spell to terminate by~$T$, so the observed
hazard \emph{diverges} at the truncation point.  A rising hazard is positive
duration dependence; positive duration dependence is under-dispersion; and by
\eqref{eq:sigmae} under-dispersion is $\sigma^2_e < 0$.  The bias is imposed
by construction, independently of the true data-generating process.

\section{Exact result for the exponential case}

The exponential is the natural benchmark because it sits exactly at
$\sigma^2_e = 0$: it is the case the decomposition treats as ``no duration
dependence''.  We show truncation moves it strictly off that point, in one
direction only.

\begin{proposition}
Let $Y \sim \mathrm{Exp}$ with mean normalised to $1$, and let
$\theta = T/E[Y] > 0$ be the truncation point in units of the mean.  Then
\[
  \operatorname{Var}(Y \mid Y \le T) \;-\; E[Y \mid Y \le T]^2 \;<\; 0
  \qquad\text{for every } \theta > 0 .
\]
\end{proposition}

\begin{proof}
Write $p = 1 - e^{-\theta}$ for the retention probability and set
\[
  a \;=\; \frac{\theta e^{-\theta}}{p} \;=\; \frac{\theta}{e^{\theta}-1}.
\]
Direct integration of the truncated density on $[0,T]$ gives
\[
  E[Y \mid Y \le T] = 1 - a,
  \qquad
  E[Y^2 \mid Y \le T] = 2 - a(2+\theta),
\]
hence
\[
  \operatorname{Var}(Y \mid Y\le T)
   = 2 - a(2+\theta) - (1-a)^2
   = 1 - a\theta - a^2 .
\]
Therefore
\begin{equation}
  \operatorname{Var} - E[\,\cdot\,]^2
   = \bigl(1 - a\theta - a^2\bigr) - \bigl(1 - a\bigr)^2
   = a\bigl(2 - \theta - 2a\bigr),
  \label{eq:gap}
\end{equation}
and since $a > 0$ the sign is that of
\[
  g(\theta) \;=\; 2 - \theta - \frac{2\theta}{e^{\theta}-1}.
\]
For $\theta \ge 2$ the first two terms are non-positive and the third is
strictly negative, so $g(\theta) < 0$.  For $0 < \theta < 2$ it suffices to
show $(2-\theta)(e^{\theta}-1) < 2\theta$.  Define
\[
  h(\theta) = 2\theta - (2-\theta)\bigl(e^{\theta}-1\bigr).
\]
Then $h(0) = 0$ and
\[
  h'(\theta) = 1 + e^{\theta}(\theta - 1),
  \qquad h'(0) = 0,
  \qquad h''(\theta) = \theta e^{\theta} > 0
  \quad (\theta > 0).
\]
So $h'$ is strictly increasing from $h'(0)=0$, giving $h' > 0$ on
$(0,\infty)$; hence $h$ is strictly increasing from $h(0)=0$, giving
$h > 0$.  Thus $g(\theta) < 0$ on $(0,2)$ as well, and therefore for all
$\theta > 0$.
\end{proof}

\paragraph{Consequences.}
Two features of \eqref{eq:gap} are worth stating explicitly.

\begin{enumerate}[leftmargin=*]
\item \textbf{There is no safe truncation point.}  The bias is strictly
negative for every finite $\theta$; it vanishes only in the limit
$\theta \to \infty$, at rate $O(\theta^2 e^{-\theta})$.  A longer window
attenuates the problem but never eliminates it.

\item \textbf{Slow types are hit hardest.}  The bias is largest in relative
terms precisely where $\theta$ is small, i.e.\ for individuals whose mean
duration is long relative to the window.  These are also the individuals most
likely to be dropped entirely, so truncation simultaneously (i)~compresses
the dispersion of those who survive and (ii)~selects the sample towards fast
types.  The first effect depresses $\sigma^2_e$; the second inflates the
measured heterogeneity share $\sigma^2_b$.
\end{enumerate}

\section{Numerical verification}

The simulation in \texttt{Example Plots.ipynb} draws $N = 100{,}000$
individuals with two spells each from a two-type exponential mixture (means
$0.3^{-1}$ and $0.03^{-1}$ months, equal shares), then drops an individual if
\emph{either} spell exceeds the limit.  At a limit of $T = 23$ (the $75$th
percentile), the proposition predicts:

\begin{center}
\begin{tabular}{lccc}
\toprule
Type & $\theta$ & Retention $p^2$ & $\sigma^2_i - \mu_i^2$ \\
\midrule
Fast (mean $3.33$)  & $6.90$ & $0.999$ & $-0.38$ \\
Slow (mean $33.33$) & $0.69$ & $0.248$ & $-60.8$ \\
\bottomrule
\end{tabular}
\end{center}

\noindent
Aggregating over the retained sample gives predicted
$N = 62{,}370$, $\sigma^2_e = -12.4$, $\sigma^2_c = 29.5$, against simulated
values of $62{,}356$, $-12.25$ and $29.4$ respectively --- agreement to
within Monte Carlo error.  The true $\sigma^2_e$ in the untruncated sample is
$+1.29$ (share $0.002$, i.e.\ correctly detecting no duration dependence);
after truncation the estimated share is $-0.50$.

The same experiment run on a Weibull sample with genuine negative duration
dependence ($k = 0.44$, true share $0.87$) yields truncated shares of $0.76$,
$0.53$ and $0.37$ at the $99$th, $90$th and $75$th percentile limits: the
bias is in the same direction and does not require the true process to be
exponential.  Whether it flips the \emph{sign} of the estimate depends on how
strong the true duration dependence is relative to the severity of
truncation.

\section{Empirical signature in our estimates}

The proposition implies a testable comparative static: holding the population
fixed, $\sigma^2_e$ should \emph{increase} as the observation window
lengthens, because $\theta$ rises for every individual.  Our own time-window
results move in exactly that direction in both countries.

\begin{center}
\begin{tabular}{lcccc}
\toprule
Spain, window & $10$yr & $16$yr & $22$yr & $29$yr \\
\midrule
$\sigma^2_e$ & $-0.049$ & $+0.016$ & $-0.002$ & $+0.103$ \\
\bottomrule
\end{tabular}
\qquad
\begin{tabular}{lccc}
\toprule
Germany, window & $5$yr & $7$yr & $10$yr \\
\midrule
$\sigma^2_e$ & $0.552$ & $0.643$ & $0.694$ \\
\bottomrule
\end{tabular}
\end{center}

\noindent
Germany is monotone across all three windows; Spain is monotone apart from a
small non-monotonicity at $22$ years.  In both cases a shorter window --- more
severe truncation --- produces a more negative (or less positive) $\sigma^2_e$,
as the mechanism predicts.

This has a direct bearing on interpretation.  Spain's baseline estimate of
mildly positive duration dependence ($\sigma^2_e = -0.049$ over the $10$-year
window) is of the same order as the truncation bias implied by the window
length, and reverses sign once the window is extended.  It should not be read
as evidence of a genuinely rising hazard.  Germany's negative duration
dependence, by contrast, is large, robust across windows, and \emph{opposite}
in sign to the bias --- so if anything it is understated by truncation.

\end{document}
